% Sección III: Calidad, preparación e integración de factores externos
% evidencia:E04,E05,E10,E26

\section{Calidad, preparación e integración de factores externos}

La confiabilidad predictiva y censal exige resolver anomalías administrativas previas en SQL y Python~\cite{mediflow2026calidad,mediflow2026informe}.

\subsection{Auditoría archivada de calidad y consistencia temporal}

La auditoría archivada del 28 de octubre de 2026~\cite{mediflow2026calidad} registró 0,00\,\% de nulos en columnas clave de \texttt{person} (\num{912205} filas), \texttt{visit\_occurrence} (\num{1706001}) y \texttt{visit\_detail} (\num{3281613}). Aunque las claves primarias no se repiten, agrupar por clave de evento reveló duplicidad administrativa en visitas principales y residual en etapas, junto con edades imposibles que distorsionaban el grupo pediátrico.

Para verificar si la anonimización desplazó fechas respecto del calendario colombiano, se midió el volumen relativo en festivos entre semana frente a tres semanas vecinas desplazando la referencia entre $-21$ y $+21$ días~\cite{mediflow2026calidad}. En Rehabilitación Ambulatoria (cierre no laboral), el mínimo global se produjo exactamente en desfase 0 días: volumen relativo de 0,0032 (0,32\,\%) en \texttt{visit\_occurrence} y 0,0033 (0,33\,\%) en \texttt{visit\_detail} (márgenes de 0,50 y 0,49 sobre el segundo menor desplazamiento), con caídas a 0,25\,\% y 0,21\,\% en lunes y martes de Carnaval. Urgencias mantiene operación continua (73,95\,\% en festivos). Esta prueba habilita variables calendáricas reales sin certificar exactitud horaria individual; todas las cifras provienen de registros archivados sin reejecución SQL.

\begin{table*}[t]
\caption{Reglas de calidad de datos, hallazgos archivados e implementación SQL vigente}
\label{tab:reglas_calidad}
\centering
\footnotesize
\begin{tabularx}{\textwidth}{@{}L{0.17\textwidth} L{0.31\textwidth} L{0.25\textwidth} Y@{}}
\toprule
Dimensión & Hallazgo archivado y población & Implementación vigente & Efecto o límite \\
\midrule
Nulos críticos & 0,00\,\% en claves y marcas de \texttt{person}, \texttt{visit\_occurrence} y \texttt{visit\_detail} & No requiere imputación de fechas & Completitud condicionada al extracto \\
Duplicados visitas & \num{228259} copias sobrantes (13,38\,\%) sobre \num{1706001} filas & \texttt{COUNT(DISTINCT (person\_id, concept, inicio))} en SQL \texttt{04}, \texttt{11}, \texttt{13} & Reduce llegadas totales a \num{1477742} ($-13{,}4\,\%$) \\
Duplicados etapas & \num{14081} copias sobrantes (0,43\,\%) sobre \num{3281613} etapas & Conteo agregado por etapa; tasa distinta a visitas & No extrapolar 0,43\,\% como tasa de visitas \\
Edades imposibles & 43 visitas (22 pacientes, 0,0025\,\%): edad $<0$ o $>110$ años & Etiqueta \texttt{invalida} en SQL \texttt{10} y \texttt{13} & Evita inflar el grupo pediátrico \texttt{0-17} \\
Etapas inválidas y huérfanas & 0\,\% sin fin o fin $<$ inicio; 117 huérfanas globales (91 en \texttt{9201}) & SQL \texttt{14}--\texttt{16}, \texttt{18} excluyen fechas inválidas; \texttt{SQL16} usa \texttt{LEFT JOIN} & Conserva huérfanas; contraste en recorte $p_{99{,}9}$ \\
Regularización & Ausencia de filas en horas o días sin movimientos & \texttt{grilla\_regular} rellena con 0 en \texttt{agregados.py} & Evita sesgar al alza promedios horarios/diarios \\
\bottomrule
\end{tabularx}
\ieeenota{Síntesis contrastada entre~\cite{mediflow2026calidad},~\cite{mediflow2026informe} (sec.~3) y los SQL \texttt{04}, \texttt{13}--\texttt{16} y \texttt{18} en \texttt{src/hun/data/sql/}.}
\end{table*}

\subsection{Depuración estadística de estancias extremas (\texorpdfstring{$p_{99{,}9}$}{p99.9}) y regla sobre huérfanas}

La auditoría sin recorte (\nolinkurl{23_duracion_etapas_percentiles.sql}) halló máximos crudos de \num{28070,3}\,h en Tratamiento Urgencias, \num{20730,3}\,h en Urgencias y \num{11213,6}\,h en Hospitalización~\cite{mediflow2026saturacion,mediflow2026informe}. Para evitar inflar la ocupación acumulada por cierres abiertos, las consultas \texttt{14}, \texttt{15} y \texttt{16} filtran por el percentil $p_{99{,}9}$ de cada etapa: \textbf{\num{1176,3}\,h} en Hospitalización (\num{1176,257222}\,h), \textbf{\num{283,7}\,h} en Tratamiento Urgencias, \textbf{\num{17,5}\,h} en Urgencias y \textbf{0,0\,h} en Rehabilitación~\cite{mediflow2026saturacion}.

En Hospitalización (\texttt{9201}), sobre \num{108601} etapas válidas se excluyen 107 superiores a \num{1176,3}\,h (todas con visita padre) y se conservan \textbf{\num{108494} etapas}, incluidas las 91 huérfanas ($<720$\,h)~\cite{mediflow2026informe}. Este filtro reduce $p_{50}/p_{75}/p_{90}$ poscalentamiento de 288 / 306 / 319 (máx. 378) a \textbf{283 / 297 / 309} (máx. \textbf{359}) en Hospitalización, y de 219 / 260 / 300 a \textbf{201 / 239 / 272} en Tratamiento Urgencias~\cite{mediflow2026informe}. Es un recorte estadístico agregado, no una adjudicación clínica. Asimismo, aunque pasajes previos de~\cite{mediflow2026calidad} recomendaban excluir huérfanas, la implementación vigente en \texttt{16\_duracion\_etapas\_diaria.sql} usa \texttt{LEFT JOIN visit\_occurrence} al igual que \texttt{SQL14} y \texttt{SQL15}, fijando la cohorte común de \num{108494} etapas (exigir visita padre reduciría la población a \num{108403}).

\subsection{Integración externa y protocolo de cuatro filtros}

El módulo \texttt{src/external/} construye \texttt{data/externos.parquet} integrando calendario colombiano (librería \texttt{holidays} con festivos, puentes y día posfestivo, Carnaval, Semana Santa, vacaciones y quincenas/primas), clima de Open-Meteo (Soledad) y SIVIGILA semanal (dengue e IRA grave con rezagos anti-fuga \texttt{lag2s} y \texttt{lag3s}). El control \texttt{data/externos\_fuentes.json} (27 de octubre de 2026) marca estado \texttt{ok} para calendario, clima (\texttt{2018-01-01} a \texttt{2026-10-21}) y epidemiología, sin implicar por sí solo aporte predictivo~\cite{mediflow2026informe}.

\begin{table*}[t]
\caption{Resultados archivados del protocolo de evaluación empírica de factores externos}
\label{tab:factores_externos}
\centering
\footnotesize
\begin{tabularx}{\textwidth}{@{}L{0.23\textwidth} l r r r r Y@{}}
\toprule
Grupo evaluado & Decisión & Mediana & Placebo & Pliegues & $p$-signo & Motivo registrado \\
\midrule
\texttt{festivos\_y\_puentes} & Acepta & $+33{,}80\,\%$ & $-1{,}80\,\%$ & 88\,\% & 0,000 & Cumple los 4 criterios \\
\texttt{carnaval\_y\_semana\_santa} & Acepta & $+3{,}49\,\%$ & $+1{,}26\,\%$ & 73\,\% & 0,014 & Cumple los 4 criterios \\
\texttt{vacaciones\_escolares} & Descarta & $+4{,}41\,\%$ & $-0{,}13\,\%$ & 65\,\% & 0,084 & $< 70\,\%$ de pliegues y $p > 0{,}05$ \\
\texttt{quincena\_y\_primas} & Descarta & $+1{,}17\,\%$ & $+0{,}79\,\%$ & 58\,\% & 0,279 & Placebo $> 0{,}5\times$, $< 70\,\%$ pliegues y $p > 0{,}05$ \\
\bottomrule
\end{tabularx}
\ieeenota{Fuente:~\cite{mediflow2026externos} (18 de noviembre de 2026, 26 pliegues \textit{walk-forward}). Criterios: mediana $\ge 1\,\%$, mejora en $\ge 70\,\%$ de pliegues, prueba unilateral de signo $p \le 0{,}05$ y placebo $< 0{,}5\times$ mejora real. \texttt{clima} y \texttt{epidemiologia} fueron descartados en una corrida inicial anterior~\cite{mediflow2026informe} y no tienen cifras en la tabla del 18 de noviembre. Las cifras de \texttt{docs/protocolo\_factores\_externos.md} ($+29{,}6\,\%$, placebo $-0{,}69\,\%$) corresponden a una corrida previa.}
\end{table*}

Cada grupo externo enfrenta cuatro filtros consecutivos sobre los 26 pliegues~\cite{mediflow2026externos}: ablación frente a base calendárica (día, mes, armónicos anuales), control de estacionalidad, prueba placebo desplazada y estabilidad/prueba de signo. Como resume la Tabla~\ref{tab:factores_externos}, solo \texttt{festivos\_y\_puentes} (6 variables) y \texttt{carnaval\_y\_semana\_santa} (5 variables) son aceptados en \texttt{lgbm\_final}.
